\begin{table}[htbp]
\centering
\footnotesize
\setlength{\tabcolsep}{3pt}
\caption{截距分解与自由幂律指数 $\gamma$，附双向自助法 95\% 置信区间。}
\label{tab:s2}
\begin{tabular}{lccccccc}
\toprule
设定 & $a$ 总量 & \begin{tabular}[t]{@{}c@{}}去除周效应\\后的 $a$\end{tabular} & \begin{tabular}[t]{@{}c@{}}去除周效应与\\州-季效应后的 $a$\end{tabular} & 周占比 & $\gamma$ & \begin{tabular}[t]{@{}c@{}}非收敛抽样\\占比（\%）\end{tabular} & \begin{tabular}[t]{@{}c@{}}仿射局部\\斜率均值\end{tabular} \\
\midrule
$R_A$，规则窗口（$S=4$） & \begin{tabular}[t]{@{}c@{}}0.0150\\ \relax[0.0100, 0.0205]\end{tabular} & \begin{tabular}[t]{@{}c@{}}0.0088\\ \relax[0.0066, 0.0122]\end{tabular} & \begin{tabular}[t]{@{}c@{}}0.0083\\ \relax[0.0066, 0.0137]\end{tabular} & \begin{tabular}[t]{@{}c@{}}0.41\\ \relax[0.12, 0.57]\end{tabular} & \begin{tabular}[t]{@{}c@{}}0.38\\ \relax[0.29, 0.51]\end{tabular} & 1.3 & $-$0.44 \\
$G$，规则窗口（$S=4$） & \begin{tabular}[t]{@{}c@{}}0.0840\\ \relax[0.0515, 0.1199]\end{tabular} & \begin{tabular}[t]{@{}c@{}}0.0415\\ \relax[0.0304, 0.0569]\end{tabular} & \begin{tabular}[t]{@{}c@{}}0.0394\\ \relax[0.0305, 0.0602]\end{tabular} & \begin{tabular}[t]{@{}c@{}}0.51\\ \relax[0.21, 0.65]\end{tabular} & \begin{tabular}[t]{@{}c@{}}0.35\\ \relax[0.21, 0.48]\end{tabular} & 0.3 & $-$0.38 \\
$R_A$，日历窗口（$S=3$） & \begin{tabular}[t]{@{}c@{}}0.0202\\ \relax[0.0104, 0.0307]\end{tabular} & \begin{tabular}[t]{@{}c@{}}0.0091\\ \relax[0.0068, 0.0126]\end{tabular} & \begin{tabular}[t]{@{}c@{}}0.0078\\ \relax[0.0056, 0.0145]\end{tabular} & \begin{tabular}[t]{@{}c@{}}0.55\\ \relax[0.11, 0.71]\end{tabular} & \begin{tabular}[t]{@{}c@{}}0.27\\（不报告）\end{tabular} & 14.4 & $-$0.33 \\
$G$，日历窗口（$S=3$） & \begin{tabular}[t]{@{}c@{}}0.0991\\ \relax[0.0608, 0.1356]\end{tabular} & \begin{tabular}[t]{@{}c@{}}0.0473\\ \relax[0.0353, 0.0652]\end{tabular} & \begin{tabular}[t]{@{}c@{}}0.0450\\ \relax[0.0326, 0.0901]\end{tabular} & \begin{tabular}[t]{@{}c@{}}0.52\\ \relax[0.24, 0.65]\end{tabular} & \begin{tabular}[t]{@{}c@{}}0.27\\（不报告）\end{tabular} & 13.8 & $-$0.33 \\
\bottomrule
\end{tabular}
\par\smallskip
\begin{minipage}{0.95\linewidth}\scriptsize 中心化层级：$a$ 总量按规格合并窗口均值；去除周效应时按同周辖区均值中心化，并乘 $n_w/(n_w-1)$ 校正自由度；州-季效应以周与州-季加性预测值（交替投影）中心化。周占比 $=1-a_{\text{去周}}/a_{\text{总量}}$，区间取自与分解相同的自助重抽；去除周效应给出统计上的方差划分，不作因果解释。$\gamma$ 为 $E(L)=cM^{-\gamma}$ 的全局拟合。非收敛抽样指 $\gamma>2$ 或 $\gamma<0$ 的自助重复；规则窗口中这类抽样不足 2\%，却使自助标准差膨胀十余倍。因此基于标准差的 Wald 检验与百分位区间对 $\gamma=1$ 的结论不一致，本表不报告检验 $p$ 值，百分位区间仅作描述。日历窗口的 $\gamma$ 自助分布严重重尾，区间不报告，亦不作推断依据。仿射局部斜率均值为 $-1/(1+M/(b/a))$ 在观测样本上的平均，仅为几何描述，与 $\gamma$ 没有相等关系。区间计算同表~\ref{tab:s1}。\end{minipage}
\end{table}
